% ===================================================================
%  Αρχείο: 15347.tex (Fragment)
%  Αυτόματη μετατροπή μέσω doc2latex.py
% ===================================================================

ΘΕΜΑ 4

Δίνεται η συνάρτηση $f(x) = 2\sigma\upsilon\nu^2(\pi - x) - 3\eta\mu\left(\frac{\pi}{2} + x\right) + \alpha$, με $\alpha \in \mathbb{R}$.

\begin{alist}
\item Να δείξετε ότι $f(x) = 2\sigma\upsilon\nu^2 x - 3\sigma\upsilon\nu x + \alpha$.

\monades{8}

\item Να εξετάσετε αν η συνάρτηση $f$ είναι άρτια ή περιττή.

\monades{5} γ) Να βρείτε το $\alpha$ αν είναι γνωστό ότι η γραφική παράσταση της $f$ διέρχεται από το σημείο $M\left(\frac{\pi}{3}, 1\right)$.

\monades{5} δ) Για $\alpha = 2$ και $g(x) = 2\eta\mu^2 x + 9\sigma\upsilon\nu x - 9$, να εξετάσετε (αν υπάρχουν) κοινά σημεία των γραφικών παραστάσεων των συναρτήσεων $f$ και $g$.

\monades{7}
\end{alist}
